% GENERATED FILE. Edit canonical lessons, then run npm run generate:books.
% canonical-source-hash: fnv1a64:486021401d95a282
% canonical-lessons: ES-C566-anecdota, ES-C566-aventura, ES-C566-hazana, ES-C566-record, ES-C566-novedad

\chapter{Una aventura}
\label{ch:es-una-aventura}
\hlchaptermodality{566}

\begin{chapteropening}
\textbf{By the end of this chapter:} \emph{I can tell an anecdote about an adventure, a great feat, a record or some news.}

Complete the last lesson of chapter 566: I can tell an anecdote about an adventure, a great feat, a record or some news.
\end{chapteropening}

\section[la anécdota]{la anécdota — the anecdote}

\label{lesson:ES-C566-anecdota}



\begin{quote}
Before the new one: say the Spanish for the episode, then the Spanish for the ending.
\end{quote}

\subsection*{You'll want to know: la anécdota}
\textbf{la anécdota} — \textquotedblleft{}the anecdote\textquotedblright{}.

\begin{quote}
Me contó una anécdota muy graciosa.
\end{quote}

\begin{quote}
\emph{He told me a very funny anecdote.}
\end{quote}

\begin{cousinweb}
\emph{anécdota} was borrowed from French \emph{anecdote}, which comes from Greek \emph{anékdota} 'unpublished things' (EtymDB), built on \emph{an-} 'not' and \emph{ékdotos} 'given out, published'.
\end{cousinweb}

\begin{grammarlens}[title={What you've built}]
One more word for telling, explaining and describing at B1.
\end{grammarlens}

\subsection*{Your turn}
\begin{itemize}
\raggedright
  \item \emph{Say it:} \emph{la anécdota}
  \item \emph{Say it:} \emph{la anécdota}, once more
  \item \emph{Say it:} a sentence of your own with \emph{la anécdota}
  \item \emph{Recall:} say the Spanish for the episode, then the Spanish for the ending, then say \emph{la anécdota} again
\end{itemize}

\begin{tcolorbox}[breakable,colback=teal!4,colframe=teal!35!black,title={Before you move on}]
What does la anécdota mean? (\textquotedblleft{}The anecdote\textquotedblright{}.) Say it once more.
\end{tcolorbox}
\section[la aventura]{la aventura — the adventure; the affair}

\label{lesson:ES-C566-aventura}



\begin{quote}
Before the new one: say the Spanish for the ending, then the Spanish for the anecdote.
\end{quote}

\subsection*{You'll want to know: la aventura}
\textbf{la aventura} — \textquotedblleft{}the adventure; the affair\textquotedblright{}.

\begin{quote}
El viaje fue toda una aventura.
\end{quote}

\begin{quote}
\emph{The trip was quite an adventure.}
\end{quote}

It can also mean a love affair.

\begin{cousinweb}
\emph{aventura} comes from spoken Latin \emph{*adventūra} 'the things that are going to happen' (EtymDB), from \emph{advenīre} 'to arrive', built on \emph{venīre} 'to come'.
\end{cousinweb}

\begin{grammarlens}[title={What you've built}]
One more word for telling, explaining and describing at B1.
\end{grammarlens}

\subsection*{Your turn}
\begin{itemize}
\raggedright
  \item \emph{Say it:} \emph{la aventura}
  \item \emph{Say it:} \emph{la aventura}, once more
  \item \emph{Say it:} a sentence of your own with \emph{la aventura}
  \item \emph{Recall:} say the Spanish for the ending, then the Spanish for the anecdote, then say \emph{la aventura} again
\end{itemize}

\begin{tcolorbox}[breakable,colback=teal!4,colframe=teal!35!black,title={Before you move on}]
What does la aventura mean? (\textquotedblleft{}The adventure; the affair\textquotedblright{}.) Say it once more.
\end{tcolorbox}
\section[la hazaña]{la hazaña — the feat, the great deed}

\label{lesson:ES-C566-hazana}



\begin{quote}
Before the new one: say the Spanish for the anecdote, then the Spanish for the adventure.
\end{quote}

\subsection*{You'll want to know: la hazaña}
\textbf{la hazaña} — \textquotedblleft{}the feat, the great deed\textquotedblright{}.

\begin{quote}
Cruzar el país a pie fue una hazaña.
\end{quote}

\begin{quote}
\emph{Crossing the country on foot was a great feat.}
\end{quote}

\begin{cousinweb}
\emph{hazaña} comes from Andalusi Arabic \emph{ḥasána} 'a good deed', from Arabic \emph{ḥasana} (EtymDB). The good deed grew into a heroic one.
\end{cousinweb}

\begin{grammarlens}[title={What you've built}]
One more word for telling, explaining and describing at B1.
\end{grammarlens}

\subsection*{Your turn}
\begin{itemize}
\raggedright
  \item \emph{Say it:} \emph{la hazaña}
  \item \emph{Say it:} \emph{la hazaña}, once more
  \item \emph{Say it:} a sentence of your own with \emph{la hazaña}
  \item \emph{Recall:} say the Spanish for the anecdote, then the Spanish for the adventure, then say \emph{la hazaña} again
\end{itemize}

\begin{tcolorbox}[breakable,colback=teal!4,colframe=teal!35!black,title={Before you move on}]
What does la hazaña mean? (\textquotedblleft{}The feat, the great deed\textquotedblright{}.) Say it once more.
\end{tcolorbox}
\section[el récord]{el récord — the record (best result)}

\label{lesson:ES-C566-record}



\begin{quote}
Before the new one: say the Spanish for the adventure, then the Spanish for the feat.
\end{quote}

\subsection*{You'll want to know: el récord}
\textbf{el récord} — \textquotedblleft{}the record (best result)\textquotedblright{}.

\begin{quote}
Tiene el récord del mundo.
\end{quote}

\begin{quote}
\emph{He holds the world record.}
\end{quote}

\emph{Tener el récord} — to hold the record. A record in writing is \emph{un registro}.

\begin{cousinweb}
\emph{récord} was borrowed from English \emph{record} (EtymDB), which came from Old French \emph{recorder} 'to recall', from Latin \emph{recordārī} 'to call to mind'.
\end{cousinweb}

\begin{grammarlens}[title={What you've built}]
One more word for telling, explaining and describing at B1.
\end{grammarlens}

\subsection*{Your turn}
\begin{itemize}
\raggedright
  \item \emph{Say it:} \emph{el récord}
  \item \emph{Say it:} \emph{el récord}, once more
  \item \emph{Say it:} a sentence of your own with \emph{el récord}
  \item \emph{Recall:} say the Spanish for the adventure, then the Spanish for the feat, then say \emph{el récord} again
\end{itemize}

\begin{tcolorbox}[breakable,colback=teal!4,colframe=teal!35!black,title={Before you move on}]
What does el récord mean? (\textquotedblleft{}The record (best result)\textquotedblright{}.) Say it once more.
\end{tcolorbox}
\section[la novedad]{la novedad — the novelty; the news, the change}

\label{lesson:ES-C566-novedad}



\begin{quote}
Before the new one: say the Spanish for the feat, then the Spanish for the record.
\end{quote}

\subsection*{You'll want to know: la novedad}
\textbf{la novedad} — \textquotedblleft{}the novelty; the news, the change\textquotedblright{}.

\begin{quote}
¿Hay alguna novedad sobre el caso?
\end{quote}

\begin{quote}
\emph{Is there any news about the case?}
\end{quote}

\emph{Sin novedad} — nothing new to report.

\begin{cousinweb}
\emph{novedad} is inherited from Latin \emph{novitās}, \emph{novitātem} 'newness' (EtymDB), built on \emph{novus} 'new'.
\end{cousinweb}

\begin{grammarlens}[title={What you've built}]
One more word for telling, explaining and describing at B1.
\end{grammarlens}

\subsection*{Your turn}
\begin{itemize}
\raggedright
  \item \emph{Say it:} \emph{la novedad}
  \item \emph{Say it:} \emph{la novedad}, once more
  \item \emph{Say it:} a sentence of your own with \emph{la novedad}
  \item \emph{Recall:} say the Spanish for the feat, then the Spanish for the record, then say \emph{la novedad} again
\end{itemize}

\begin{tcolorbox}[breakable,colback=teal!4,colframe=teal!35!black,title={Before you move on}]
What does la novedad mean? (\textquotedblleft{}The novelty; the news, the change\textquotedblright{}.) Say it once more.
\end{tcolorbox}
